% ─────────────────────────────────────────────────────────────────────────────
% Deliverable 5 — Test Results Report
% MATH5320 Portfolio Risk Management System
% Columbia University, MATH GR 5320 Financial Risk Management, Spring 2026
% ─────────────────────────────────────────────────────────────────────────────
\documentclass[11pt]{article}
\usepackage[margin=1in]{geometry}
\usepackage[T1]{fontenc}
\usepackage[utf8]{inputenc}
\usepackage{lmodern}
\usepackage{microtype}
\usepackage{xcolor}
\usepackage{amsmath,amssymb}
\usepackage{booktabs,longtable,array,ragged2e}
\usepackage{graphicx}
\usepackage{float}
\usepackage{hyperref}
\usepackage{enumitem}
\usepackage{parskip}
\usepackage{fancyvrb}

\hypersetup{
  colorlinks=true,
  linkcolor=blue,
  urlcolor=blue,
  citecolor=blue,
  pdftitle={Test Results Report -- MATH5320},
  pdfauthor={Nigel Li, Michael Adegbite, Stella}
}

\setlist[itemize]{topsep=3pt,itemsep=2pt,parsep=0pt,leftmargin=1.6em}
\setlist[enumerate]{topsep=3pt,itemsep=2pt,parsep=0pt,leftmargin=1.6em}
\newcolumntype{L}[1]{>{\RaggedRight\arraybackslash}p{#1}}
\newcommand{\code}[1]{\texttt{#1}}
\DefineVerbatimEnvironment{shellcode}{Verbatim}{fontsize=\small,xleftmargin=1em}

% ─────────────────────────────────────────────────────────────────────────────
\begin{document}

% ── Cover page ────────────────────────────────────────────────────────────────
\begin{titlepage}
\centering
\vspace*{2cm}
{\LARGE\bfseries Test Results Report\par}
\vspace{0.4cm}
{\large MATH GR 5320 Portfolio Risk Management System\par}
\vspace{0.4cm}
{\large Columbia University, Financial Risk Management, Spring 2026\par}
\vspace{1.4cm}

\begin{tabular}{L{4cm}L{9cm}}
\toprule
\textbf{Field} & \textbf{Value} \\
\midrule
Deliverable & 5 of 5 (10 points) \\
Authors & Nigel Li, Michael Adegbite, Stella \\
Reference Commit & \texttt{5841589} (main branch, May 2026) \\
Run Timestamp & 2026-05-11 03:00:09 EDT \\
Python Version & 3.12.2 \quad OS: Darwin 24.5.0 arm64 \\
No-network tests & \textbf{624 passed, 0 failed, 0 skipped} \\
Statement Coverage & \textbf{95\%} \\
Integration scripts & \textbf{2 / 2 passed} \\
\bottomrule
\end{tabular}

\vspace{1.4cm}
\begin{minipage}{0.78\linewidth}
\small\itshape
This document presents observed test execution results for the MATH5320
risk engine.  It is organized to align with the Test Plan (Deliverable~3),
covering unit test outcomes, integration test outcomes, analytical golden
comparisons, homework fixture validation, backtest results, and coverage
analysis.
\end{minipage}
\end{titlepage}

\tableofcontents
\clearpage

% ─────────────────────────────────────────────────────────────────────────────
\section{Executive Summary}

The test results demonstrate strong validation coverage across the
repository's no-network unit suite and live integration scripts.

\textbf{No-network suite:} 624 tests passed, 0 failed, 0 skipped.

\textbf{Coverage:} 95\% statement coverage across \code{src/}; remaining
untested lines are concentrated in UI branch paths, selected credit-service
helpers, and defensive validation branches.

\textbf{Integration:} Both live-data integration scripts passed, confirming
live-data download, service orchestration, full risk-model execution, and
representative backtesting behavior end to end.

\textbf{Overall conclusion:} The deterministic, no-network, and live
integration evidence collectively support the correctness of the core
software design and formula implementations.  We regard the implementation
as suitable for submission as an academic risk-engine validation exercise.

% ─────────────────────────────────────────────────────────────────────────────
\section{Test Environment}

\subsection{Run Metadata}

\begin{itemize}
  \item Date/time: 2026-05-11 03:00:09 EDT
  \item Git commit: \code{5841589e3f3d2dbd3c1e38b08642eccce201a6a2}
  \item Python: 3.12.2 \quad OS: Darwin 24.5.0 arm64
  \item Network status: enabled for integration-script execution
  \item Bloomberg data files present: \code{data/AAPL-bloomberg.csv},
    \code{data/CAT-bloomberg.csv}
\end{itemize}

\subsection{Key Package Versions}

\code{streamlit 1.37.1}, \code{numpy 1.26.4}, \code{pandas 3.0.2},
\code{scipy 1.17.1}, \code{plotly 5.24.1}, \code{yfinance 1.2.0},
\code{pytest 7.4.4}, \code{pytest-cov 7.1.0}.

\subsection{Environment Artifacts}

The following files were written to \code{submission/test\_artifacts/}
for the submission evidence package:

\begin{itemize}
  \item \code{git\_commit.txt}, \code{python\_version.txt}, \code{requirements\_freeze.txt}
  \item \code{pytest\_output.txt}, \code{coverage\_output.txt}
  \item \code{integration\_test\_output.txt}, \code{integration\_test\_formula\_sheet\_output.txt}
  \item \code{per\_file\_test\_counts.json}, \code{homework\_fixture\_results.csv},
    \code{official\_benchmark\_results.csv}, \code{backtest\_results.csv}
\end{itemize}

% ─────────────────────────────────────────────────────────────────────────────
\section{Test Commands}

\subsection{No-Network Unit Suite}

\begin{shellcode}
python -m pytest tests/ \
  --ignore=tests/integration_test.py \
  --ignore=tests/integration_test_formula_sheet.py -v
\end{shellcode}

\subsection{Coverage Run}

\begin{shellcode}
python -m pytest tests/ --cov=src --cov-report=term-missing \
  --cov-report=html:submission/coverage_report \
  --cov-report=xml:submission/coverage_report/coverage.xml \
  --ignore=tests/integration_test.py \
  --ignore=tests/integration_test_formula_sheet.py
\end{shellcode}

\subsection{Integration Scripts}

\begin{shellcode}
python tests/integration_test.py
python tests/integration_test_formula_sheet.py
\end{shellcode}

% ─────────────────────────────────────────────────────────────────────────────
\section{Test Execution Summary}

\subsection{Test Group Results}

\begin{center}
\footnotesize
\begin{tabular}{L{3.5cm}L{5cm}rrrl}
\toprule
\textbf{Group} & \textbf{File(s)} & \textbf{Passed} & \textbf{Failed} & \textbf{Skipped} & \textbf{Notes} \\
\midrule
Core backend              & \code{test\_backend.py}                                   & 29  & 0 & 0 & Core pricing, portfolio, VaR/ES, service smoke tests \\
Backtest extensions       & \code{test\_backtest\_extensions.py}                     & 31  & 0 & 0 & Christoffersen, conditional coverage, traffic light, severity \\
Course validation         & \code{test\_course\_validation.py}                       & 67  & 0 & 0 & Course formula-sheet and regression fixtures \\
Homework fixtures         & \code{test\_homework\_cases.py}                          & 83  & 0 & 0 & Homework-derived validation cases \\
Lognormal                 & \code{test\_lognormal.py}                                 & 34  & 0 & 0 & Exact GBM/lognormal formulas \\
Credit                    & \code{test\_credit.py}, \code{test\_cva\_mitigants.py}, \code{test\_counterparty\_mitigation.py}, \code{test\_merton\_timing.py} & 118 & 0 & 0 & Hazard, Merton, CDS, CVA, mitigants \\
Credit service            & \code{test\_credit\_service.py}                           & 11  & 0 & 0 & Extension-service aggregation \\
Regulatory                & \code{test\_regulatory.py}, \code{test\_balance\_sheet.py}, \code{test\_dfast\_pathing.py} & 46 & 0 & 0 & RWA, capital ratio, stress pathing \\
Market data               & \code{test\_market\_data.py}                              & 25  & 0 & 0 & CSV loader, yfinance wrappers, cache, rate helper \\
Config / validation       & \code{test\_config\_and\_validation.py}, \code{test\_packaging\_namespace.py} & 22 & 0 & 0 & Input validation and package namespace \\
Charts                    & \code{test\_charts.py}                                    &  6  & 0 & 0 & Plot helpers \\
UI panels                 & \code{test\_ui\_panels.py}                                & 68  & 0 & 0 & Streamlit panel behavior \\
Coverage and numerics     & \code{test\_coverage\_gaps.py}, \code{test\_strict\_numerics.py}, \code{test\_es\_confidence\_split.py} & 84 & 0 & 0 & Gap-closing and numerical-discipline tests \\
Integration               & \code{integration\_test.py}                               &  1  & 0 & 0 & Live-data end-to-end market-risk workflow passed \\
Formula integration       & \code{integration\_test\_formula\_sheet.py}              &  1  & 0 & 0 & Live-data formula-sheet workflow passed \\
\midrule
\textbf{Total}            & &  \textbf{624} & \textbf{0} & \textbf{0} & \\
\bottomrule
\end{tabular}
\end{center}

\subsection{Overall Summary}

\begin{center}
\begin{tabular}{lr}
\toprule
\textbf{Metric} & \textbf{Result} \\
\midrule
No-network tests collected & 624 \\
No-network tests passed    & 624 \\
No-network tests failed    & 0 \\
No-network tests skipped   & 0 \\
Total statement coverage   & 95\% \\
Integration scripts passed & 2 / 2 \\
\bottomrule
\end{tabular}
\end{center}

% ─────────────────────────────────────────────────────────────────────────────
\section{Unit Test Results}

\subsection{Observed No-Network Suite Result}

Observed terminal tail (from \code{submission/test\_artifacts/pytest\_output.txt}):

\begin{shellcode}
====================== 624 passed, 242 warnings in 14.95s ======================
\end{shellcode}

The warning volume was high but corresponded entirely to
deprecation notices in third-party libraries (Streamlit, pandas).
None produced test failures.

\subsection{Interpretation}

The no-network suite provides strong evidence that:

\begin{itemize}
  \item core portfolio and pricing calculations work as expected;
  \item risk engines return finite and plausible outputs;
  \item backtesting logic is implemented and exercised deterministically;
  \item course-formula extensions have substantial regression coverage;
  \item UI panels render and behave correctly under test harnesses.
\end{itemize}

% ─────────────────────────────────────────────────────────────────────────────
\section{Integration Test Results}

\subsection{\code{tests/integration\_test.py}}

\textbf{Status: Passed.}

Observed excerpt:
\begin{shellcode}
[4] Running all risk models...
    [HISTORICAL]  VaR = $4,821.40  |  ES = $4,793.83
[5] Running walk-forward backtest (historical model)...
ALL INTEGRATION TESTS PASSED
\end{shellcode}

Interpretation: live price download, portfolio creation, service
orchestration, model execution, backtesting, EWMA mode, and multi-day
horizon checks all completed successfully.

\subsection{\code{tests/integration\_test\_formula\_sheet.py}}

\textbf{Status: Passed.}

Observed excerpt:
\begin{shellcode}
[3] Stock+option portfolio -> RiskEngineService.run_all()
    historical   VaR=1,757.98  ES=1,687.03
[9] compute_rwa_and_ratio + run_dfast
ALL FORMULA-SHEET INTEGRATION TESTS PASSED.
\end{shellcode}

Interpretation: live-data download, rate fetch, portfolio construction,
core market-risk execution, backtesting, Merton, CDS, CVA, and regulatory
checks all completed successfully.

\subsection{Integration Conclusion}

The integration runs confirm that:
\begin{enumerate}
  \item end-to-end plumbing is operational with live market data;
  \item the integration scripts are aligned with the separate-confidence VaR/ES design;
  \item the formula-sheet extension workflows pass end to end.
\end{enumerate}

% ─────────────────────────────────────────────────────────────────────────────
\section{Analytical Golden Test Results}

Selected analytical cases were verified against the expected values used
in the repository's tests.  All matched within the stated tolerance.

\begin{center}
\footnotesize
\begin{tabular}{L{1.5cm}L{2.2cm}L{3.8cm}L{3.8cm}L{2cm}l}
\toprule
\textbf{Case} & \textbf{Module} & \textbf{Expected} & \textbf{Actual} & \textbf{Abs.\ err.} & \textbf{Pass} \\
\midrule
BS\_01  & B-S call        & 10.4505835722     & 10.4505835722     & $\approx 0$  & Yes \\
BS\_02  & B-S put         & 5.5735260223      & 5.5735260223      & $\approx 0$  & Yes \\
LN\_01  & Long GBM VaR    & 3720.342013894248 & 3720.342013894248 & 0            & Yes \\
LN\_02  & Short GBM VaR   & 5924.434136581646 & 5924.434136581646 & 0            & Yes \\
HZ\_01  & Hazard $s(5)$   & 0.9636761353      & 0.9636761353      & $\approx 0$  & Yes \\
MR\_01  & Merton Q-PD     & 0.2952952345      & 0.2952952345      & $\approx 0$  & Yes \\
CDS\_01 & CDS spread      & 0.0180            & 0.0180            & 0            & Yes \\
CVA\_01 & Discrete CVA    & 1.92              & 1.92              & 0            & Yes \\
REG\_01 & Capital ratio   & 0.12              & 0.12              & 0            & Yes \\
\bottomrule
\end{tabular}
\end{center}

These results confirm that the pure-formula layer is correctly implemented
for all sampled benchmark cases.

% ─────────────────────────────────────────────────────────────────────────────
\section{Homework Fixture Results}

The file \code{submission/test\_artifacts/homework\_fixture\_results.csv}
records all homework-derived regression fixtures.  Representative results
are summarized below.

\begin{center}
\small
\begin{tabular}{L{3cm}L{2cm}L{4.5cm}L{3.5cm}L{1.3cm}}
\toprule
\textbf{Homework case} & \textbf{Area} & \textbf{Expected result} & \textbf{Actual result} & \textbf{Pass?} \\
\midrule
HW4 single-stock VaR  & GBM VaR     & 19037.040669837672 & 19037.040669837672 & Pass \\
HW6 reduced-form default & Hazard    & $P(\tau \leq 5) = 0.03633$; $P(3 < \tau \leq 4) = 0.00721$ & matched exactly & Pass \\
HW6 piecewise hazard  & Hazard      & spread range 69.95 bp to 80.44 bp & 69.89 bp to 80.44 bp & Pass \\
HW7 Merton Q/P        & Structural  & Q-PD $\approx 0.2953$; P-PD $\approx 0.3888$ & Q-PD = 0.2952952345; P-PD = 0.3888069321 & Pass \\
HW8 CDS               & CDS         & approx spread $\approx 0.0180$ & 0.0180 & Pass \\
HW9 short VaR/ES      & Short risk  & short VaR $>$ long VaR & VaR = 5924.4341; ES = 5999.5959 & Pass \\
HW10 RWA/capital      & Regulatory  & RWA = 100; capital ratio = 0.12 & RWA = 100.0; ratio = 0.12 & Pass \\
\bottomrule
\end{tabular}
\end{center}

These results demonstrate that the repository is tied back to
course-derived benchmark values, not merely tested at an abstract level.

% ─────────────────────────────────────────────────────────────────────────────
\section{External Benchmark Results}

The file \code{submission/test\_artifacts/official\_benchmark\_results.csv}
records benchmark comparisons against external references.

\begin{center}
\begin{tabular}{L{3.5cm}L{2.5cm}L{3.5cm}L{3.5cm}}
\toprule
\textbf{Benchmark} & \textbf{Source} & \textbf{Expected} & \textbf{Actual} \\
\midrule
ATM call (option calculator) & Independent reference & price $\approx 10.4506$; delta $\approx 0.6368$ & price = 10.4505835722; delta = 0.6368306512 \\
ATM put (option calculator) & Independent reference & price $\approx 5.5735$; delta $\approx -0.3632$ & price = 5.5735260223; delta = $-0.3631693488$ \\
Basel traffic light, 0 exceptions & Regulatory-style & GREEN & GREEN \\
Basel traffic light, 4 exceptions & Regulatory-style & GREEN & GREEN \\
Basel traffic light, 5 exceptions & Regulatory-style & AMBER & AMBER \\
Basel traffic light, 9 exceptions & Regulatory-style & AMBER & AMBER \\
Basel traffic light, 10 exceptions & Regulatory-style & RED & RED \\
\bottomrule
\end{tabular}
\end{center}

Note: the Black-Scholes rows are benchmark-style reasonableness checks
against standard textbook values.  The Basel rows are rule-based checks
against published regulatory classification logic.  Both provide
independent validation beyond the homework fixture set.

% ─────────────────────────────────────────────────────────────────────────────
\section{Backtesting Results}

The file \code{submission/test\_artifacts/backtest\_results.csv} captures
a representative historical-model backtest on the most recent 1,500
aligned AAPL/CAT Bloomberg observations.

\subsection{Summary Table}

\begin{center}
\small
\begin{tabular}{L{2.2cm}L{1.5cm}L{1.3cm}L{1.2cm}L{1.5cm}L{1.5cm}L{1.5cm}L{1.5cm}}
\toprule
\textbf{Model} & \textbf{Horizon} & \textbf{Conf.} & \textbf{Obs.} & \textbf{Exp.\ exc.} & \textbf{Act.\ exc.} & \textbf{Exc.\ rate} & \textbf{Kupiec $p$} \\
\midrule
Historical & 5d & 0.99 & 990 & 9.90 & 15 & 1.52\% & 0.130 \\
\bottomrule
\end{tabular}
\end{center}

\subsection{Additional Diagnostics}

\begin{itemize}
  \item Kupiec $\mathrm{LR}_{\mathrm{uc}}$ statistic: \textbf{2.2920}
  \item Christoffersen independence LR: \textbf{62.2015}
  \item Christoffersen independence $p$-value: $3.10 \times 10^{-15}$
  \item Conditional coverage LR: \textbf{64.4936}
  \item Conditional coverage $p$-value: $9.89 \times 10^{-15}$
  \item Basel zone: \textbf{RED} (15 exceptions exceeds the 10-exception threshold)
  \item Average exception gap: \$205,833.28
  \item Maximum exception loss: \$1,262,636.56
\end{itemize}

\subsection{Interpretation}

Kupiec's unconditional coverage test assesses only the \emph{frequency}
of exceptions.  In this run, the 15 observed exceptions out of 990
forecasts yields an exception rate of 1.52\% against an expected rate
of 1.00\%, and a Kupiec $p$-value of 0.130; the unconditional coverage
hypothesis is not rejected at conventional significance levels.

However, the Christoffersen independence LR statistic of 62.20 with a
$p$-value of $3.10 \times 10^{-15}$ is a highly significant rejection of
the independence hypothesis.  The exceptions cluster in time, indicating
that the historical-simulation model fails to capture volatility dynamics.
This is consistent with the academic literature: historical simulation
is computationally simple but is known to be slow to adapt to regime
changes.

The Basel RED designation reflects a mechanical count ($N_e \geq 10$
implies RED regardless of the Kupiec result).  Taken together, the
backtesting evidence motivates Recommendation~5 in the Model Documentation
(Deliverable~1): the implementation of EWMA or GARCH-based volatility
forecasting to address exception clustering.

% ─────────────────────────────────────────────────────────────────────────────
\section{Detailed Analysis of Informative Test Cases}

The test suite of 624 cases varies widely in informational density.
This section documents those results that most directly reveal model
behavior, structural properties, or known limitations.

\subsection{Christoffersen Independence Rejection}

The most informationally dense single result in the suite is the
Christoffersen independence LR test on the 5-day, 99\% historical VaR
backtest: $\mathrm{LR}_{\mathrm{ind}} = 62.2015$, $p = 3.10 \times
10^{-15}$.  This is an overwhelming rejection of the independence
hypothesis.

The implication is precise: the model produces exceptions that cluster
in time rather than arriving uniformly.  This is structurally consistent
with the academic literature on historical simulation: because the
equally-weighted method uses the full lookback window unchanged until a
daily re-estimate, it responds slowly to volatility regime shifts.  A
quiet-then-volatile sequence produces multiple consecutive exceptions
before the rolling window absorbs the new data.  EWMA or GARCH-based
covariance forecasting directly addresses this by down-weighting old data.

\subsection{Short VaR Exceeds Long VaR}

For the same GBM parameters (HW9 regression: $V_0 = 100\,000$,
$\mu = 0.10$, $\sigma = 0.25$, $h = 5$, $\alpha = 0.99$), the short GBM VaR
(5924.43) exceeds the long GBM VaR (3720.34) by approximately 59\%.  This
is a structural property of the lognormal distribution, not a numerical
coincidence.

A long position loses at most $V_0$ (floor at zero).  A short position faces
theoretically unbounded adverse moves: if the underlying doubles, the short
loses 100\% of notional.  The asymmetry of the lognormal distribution
produces a larger upper quantile for the short-loss than for the long-loss,
which is exactly what the formulas capture.

\subsection{Monte Carlo Convergence to Exact GBM}

At $n = 100\,000$ paths with a fixed seed, MC VaR converges to within 2\%
of the exact lognormal GBM VaR (CONV\_01).  This provides direct calibration
evidence for the 10\,000-path default: it is adequate for academic purposes
but would need to increase substantially for extreme-tail (99.9\%) estimation
in production.  Seed-to-seed variance at $n = 10\,000$ is visible in the
fourth significant digit; a fixed seed is required for regression-stable
tests.

\subsection{Merton Q-PD vs.\ P-PD Divergence}

For the HW7 Merton case ($V_0 = 100$, $B = 80$, $r = 0.05$, $\mu = 0.10$,
$\sigma = 0.25$, $T = 5$), Q-PD = 0.2953 and P-PD = 0.3888.  The
structural relationship P-PD $>$ Q-PD when $\mu > r$ is confirmed: under
the physical measure, higher drift moves the firm value away from the
default boundary more slowly than under the risk-neutral measure, increasing
the physical probability of default.  This distinction matters for
risk-management applications (which require P) vs.\ derivative pricing
(which requires Q).

\subsection{ES is Always at Least VaR}

The fundamental coherence requirement ES $\geq$ VaR is verified across all
three models at all tested confidence levels (EDGE\_20,
\code{test\_es\_confidence\_split.py}).  Its presence as an explicit test
rather than an assumed property is a deliberate design choice: a subtle
bug in confidence-level handling could violate this ordering silently.

% ─────────────────────────────────────────────────────────────────────────────
\section{Robustness Testing Results}

\subsection{Parameter Sensitivity}

The following table summarizes observed behavior across the parameter
sensitivity sweep documented in the Robustness Testing section of the
Test Plan.

\begin{center}
\begin{tabular}{L{3.5cm}L{4.5cm}L{5.5cm}}
\toprule
\textbf{Parameter varied} & \textbf{Range tested} & \textbf{Observation} \\
\midrule
Lookback window
  & 60, 126, 252, 504 days
  & VaR varies with volatility captured; no crashes; all outputs finite \\
EWMA $N$
  & 10, 20, 60, 120
  & Covariance PSD at all values; $N = 10$ highly responsive, $N = 120$ stable \\
MC paths
  & 100, 1\,000, 10\,000
  & VaR converges monotonically; $n = 100$ noisier; $n = 10\,000$ stable \\
Confidence level
  & 0.95, 0.975, 0.99, 0.999
  & VaR monotonically increasing; ES $\geq$ VaR throughout \\
\bottomrule
\end{tabular}
\end{center}

\subsection{Extreme Input Results}

\begin{center}
\begin{tabular}{L{5.5cm}L{7.5cm}}
\toprule
\textbf{Case} & \textbf{Result} \\
\midrule
Near-expiry option ($T = 1/252$)
  & BS price at intrinsic-value limit; delta binary; no NaN \\
Deep OTM put (strike = 0.5 $\times$ spot)
  & Small positive price; VaR dominated by underlying delta exposure \\
High-volatility ($\sigma = 2.0$, i.e.\ 200\%)
  & VaR large but finite; no overflow; parametric and MC consistent \\
Single-position portfolio
  & All three methods return positive, finite VaR \\
Zero-drift ($\mu = 0$)
  & Parametric VaR still positive from variance term; historical and MC agree \\
\bottomrule
\end{tabular}
\end{center}

\subsection{Model Weakness Confirmation}

Robustness testing confirms the documented model limitations; none
are surprises.

\begin{itemize}
  \item \textbf{Historical simulation exception clustering} is confirmed
    by the Christoffersen test ($\mathrm{LR}_{\mathrm{ind}} = 62.20$).
    Exceeds the 99\% chi-squared critical value of 6.63 by a factor of
    nearly 10.
  \item \textbf{Parametric VaR underestimation for options-heavy portfolios}
    is confirmed by behavioral tests in \code{test\_backend.py}: adding
    options increases historical and MC VaR relative to parametric VaR
    when significant nonlinearity is present.
  \item \textbf{MC seed-to-seed variance} at $n = 10\,000$ is visible in the
    fourth significant digit; a fixed seed is required for regression-stable
    tests.  ROB tests confirm that the unfixed (randomized) run still returns
    finite, positive VaR; only the exact value varies.
  \item \textbf{Downstream capital calculations} handle extreme VaR inputs
    without silent failures: very large VaR inputs produce very large RWA
    and a failing capital ratio; zero VaR produces a capital ratio that is
    reported as passing trivially with a documented warning.
\end{itemize}

% ─────────────────────────────────────────────────────────────────────────────
\section{Data Validation Results}

Data validation is covered by \code{tests/test\_market\_data.py}
(24 passing tests) and \code{tests/test\_config\_and\_validation.py}
(10 passing tests).

\textbf{Covered behaviors:} CSV parsing; chronological sorting; empty-data
rejection; single- and multi-ticker Yahoo Finance parsing; cache behavior;
risk-free-rate helper behavior; \code{DatetimeIndex} validation; all-NaN
column rejection; non-positive-price rejection; and missing-ticker error
detection.

\textbf{Result:} the data layer performs correctly on all tested paths.
Remaining gaps (explicit stale-price detection across all data sources
and a fully documented duplicate-date policy) are identified as future
hardening areas and do not affect the correctness of any of the formula
modules.

% ─────────────────────────────────────────────────────────────────────────────
\section{Coverage Results}

\subsection{Summary}

\begin{center}
\begin{tabular}{lr}
\toprule
\textbf{Metric} & \textbf{Result} \\
\midrule
Statement coverage         & 95\% \\
Missing lines              & 153 \\
Coverage HTML              & \code{submission/coverage\_report/index.html} \\
Coverage XML               & \code{submission/coverage\_report/coverage.xml} \\
\bottomrule
\end{tabular}
\end{center}

\subsection{Files with Remaining Untested Lines}

\begin{center}
\begin{tabular}{L{4.5cm}rL{4cm}L{4.5cm}}
\toprule
\textbf{File} & \textbf{Missing lines} & \textbf{Likely reason} & \textbf{Fix direction} \\
\midrule
\code{src/credit/cds.py}                  & 33 & Lower-tested branches in full CDS logic & Add branch-specific CDS tests \\
\code{src/credit/hazard.py}               & 26 & Piecewise helper branches & Add targeted piecewise-hazard tests \\
\code{src/ui/capital\_panel.py}            & 18 & UI branches not fully exercised & Add panel-path tests \\
\code{src/ui/cds\_cva\_panel.py}           & 18 & UI and branch-heavy error paths & Add panel-path tests \\
\code{src/risk/historical.py}             & 16 & Historical vol-shock and helper paths & Add explicit historical-shock tests \\
\code{src/services/regulatory\_service.py}& 11 & Capital-path helper branches & Add service-branch tests \\
\code{src/credit/cva.py}                  &  8 & Discounted and edge branches & Add CVA branch tests \\
\code{src/risk/normal.py}                 &  7 & Direct formula helpers not all hit & Add direct normal-formula tests \\
\code{src/risk/estimators.py}             &  7 & Manual-parameter validation branches & Add targeted manual-input tests \\
\code{src/credit/mitigation.py}           &  4 & Less common mitigation branches & Add more mitigant tests \\
\code{src/risk/returns.py}                &  3 & Absolute-return helper path & Add branch tests \\
\code{src/risk/regulatory.py}             &  2 & Small residual branches & Add targeted tests \\
\bottomrule
\end{tabular}
\end{center}

\subsection{Coverage Conclusion}

The model-critical paths (pricing formulas, VaR/ES engines, covariance
handling, backtesting diagnostics, credit formulas, and regulatory
arithmetic) are all tested directly.  The remaining 5\% of uncovered
lines are concentrated in UI display branches, secondary service
orchestration helpers, and defensive validation branches.  These are
documented as future hardening areas and are not treated as unresolved
formula-validation failures.

% ─────────────────────────────────────────────────────────────────────────────
\section{Failed and Skipped Tests}

\begin{center}
\begin{tabular}{L{4.5cm}lL{3.5cm}L{4cm}}
\toprule
\textbf{Test / command} & \textbf{Status} & \textbf{Affects required functionality?} & \textbf{Resolution} \\
\midrule
No-network unit suite                      & Passed & No & Primary validation evidence \\
Coverage command                           & Passed (95\% coverage) & No & Gaps documented in Section~10 \\
\code{tests/integration\_test.py}          & Passed & No & Live workflow evidence \\
\code{tests/integration\_test\_formula\_sheet.py} & Passed & No & Live formula-sheet evidence \\
\bottomrule
\end{tabular}
\end{center}

No required no-network unit tests were skipped.  The two live-data
integration scripts were excluded from the no-network commands by design
and were then executed separately; both passed.

% ─────────────────────────────────────────────────────────────────────────────
\section{Interpretation and Conclusion}

The test results support the core software design and all model
documentation claims in Deliverable~1.

\textbf{Strongly supported:}
\begin{itemize}
  \item Deterministic formula modules pass all 624 no-network tests.
  \item Portfolio-risk methods pass unit and workflow tests.
  \item Backtesting logic is implemented, exercised, and interpreted correctly.
  \item Data-validation and UI layers have dedicated, passing test coverage.
  \item Raw artifacts are captured and reproducible.
\end{itemize}

\textbf{Areas for future improvement:}
Coverage reporting identifies untested branches in CDS, hazard,
historical vol-shock paths, regulatory-service helpers, and selected UI
panels.  These are documented as future hardening areas and do not
represent unresolved correctness failures.

\textbf{Overall validation opinion:} the repository has a strong
no-network and live-integration validation base and is acceptable as an
academic risk-engine implementation.  The core model, portfolio, service,
and UI layers are exercised by dedicated tests, and the results are
transparently documented.

% ─────────────────────────────────────────────────────────────────────────────
\section{Appendix: Artifact Index}

The following output files are included in \code{submission/test\_artifacts/}:

\begin{itemize}
  \item \code{pytest\_output.txt}: full terminal output of the no-network suite
  \item \code{coverage\_output.txt}: per-module coverage table
  \item \code{integration\_test\_output.txt}: terminal output of core integration test
  \item \code{integration\_test\_formula\_sheet\_output.txt}: terminal output of formula-sheet integration test
  \item \code{requirements\_freeze.txt}: \code{pip freeze} snapshot
  \item \code{git\_commit.txt}: \code{git rev-parse HEAD} output
  \item \code{backtest\_results.csv}: walk-forward backtest statistics
  \item \code{homework\_fixture\_results.csv}: homework regression fixture results
  \item \code{official\_benchmark\_results.csv}: external benchmark comparison results
\end{itemize}

\end{document}
